\documentclass[11pt,a4paper]{article}
\usepackage{amsmath,amsxtra,amsthm,amssymb,mathrsfs,bbm}
\usepackage{hyperref}
\usepackage[margin=2cm]{geometry}
\usepackage{graphicx}
\usepackage[pdftex]{color}
\usepackage[normalem]{ulem}
\usepackage[shortlabels]{enumitem}
%\usepackage[small,nohug,noPostScript]{diagrams}
%\usepackage{stmaryrd}
\usepackage{makeidx}
\usepackage[all]{xy}
\usepackage{float}
\usepackage{tocloft}
\usepackage{tikz}
%\usetikzlibrary{positioning}
\usetikzlibrary{shapes,arrows}
\usetikzlibrary{patterns}
\usetikzlibrary{decorations.shapes}
\usetikzlibrary{decorations.pathreplacing}
\usetikzlibrary{decorations.pathmorphing}
\numberwithin{equation}{section}

\setlength{\parindent}{0mm}
\setlength{\parskip}{11pt}

%set default enumerate environment
\setlist[enumerate]{label={\rm(\arabic*)},itemsep=3pt,topsep=-6pt}
\setlist[itemize]{itemsep=3pt,topsep=-6pt}

\newenvironment{claim}[1][]
{\begin{list}{}{\setlength{\itemsep}{0pt}
	\setlength{\listparindent}{0pt}
	\setlength{\labelsep}{0.2cm}
	\setlength{\topsep}{2pt}
	\setlength{\leftmargin}{0.5cm}}

	\item[] \textbf{Claim:$\mbox{}\quad$}#1
	\item[] \textbf{Proof of Claim: }\\
}
{\end{list}}

%%% set up proof environment
\makeatletter
\renewenvironment{proof}[1][\proofname]{\par
    \pushQED{\qed}%
    \normalfont \topsep0\p@\@plus6\p@ \labelsep1em\relax
    \trivlist
    \item[\hskip\labelsep\indent
        \bfseries #1]\ignorespaces
    \mbox{}
}{
    \popQED\endtrivlist\@endpefalse
}
\makeatother

%%% Remove section numbering but retain section counting
%\makeatletter
%\def\@seccntformat#1{%
%  \expandafter\ifx\csname c@#1\endcsname\c@section\else
%  \csname the#1\endcsname\quad
%  \fi}
%\makeatother

%add dots to sections in TOC
\renewcommand{\cftsecleader}{\cftdotfill{\cftdotsep}}
%remove section number in TOC
\makeatletter
\renewcommand{\cftsecpresnum}{\begin{lrbox}{\@tempboxa}}
\renewcommand{\cftsecaftersnum}{\end{lrbox}}
\makeatother

%abovespace,belowspace,bodyfont,indent,headfont,headpunct(e.g.\newline),headspace,custom-head-spec
\newtheoremstyle{all}{10pt}{0pt}{\itshape}{0pt}
{\bfseries}{.}{5pt plus 1pt minus 1pt}{\thmname{#1}\thmnumber{ #2}\thmnote{ (#3)}}
\theoremstyle{all}

\newtheorem{theorem}{Theorem}[section]
\newtheorem{proposition}[theorem]{Proposition}
\newtheorem{lemma}[theorem]{Lemma}
\newtheorem{corollary}[theorem]{Corollary}
\newtheorem{example}[theorem]{Example}
\newtheorem*{example*}{Example}
\newtheorem{definition}[theorem]{Definition}
\newtheorem{exercise}[theorem]{Exercise}
\newtheorem*{exercise*}{Exercise}
\newtheorem{assumption}[theorem]{Assumption}
\newtheorem*{convention}{Convention}
\newtheorem*{notation}{Notation}

\theoremstyle{remark}
\newtheorem{remark}[theorem]{Remark}
\newtheorem*{remark*}{Remark}

\DeclareMathOperator{\isom}{\cong}
\DeclareMathOperator{\ifonlyif}{\Leftrightarrow}
\DeclareMathOperator{\imply}{\Rightarrow}
\DeclareMathOperator{\onto}{\twoheadrightarrow}
\DeclareMathOperator{\into}{\hookrightarrow}
\newcommand{\ses}[3]{0\rightarrow #1 \rightarrow #2 \rightarrow #3 \rightarrow 0}

\renewcommand{\P}{\mathbb{P}}
\newcommand{\Z}{\mathbb{Z}}
\newcommand{\R}{\mathbb{R}}
\newcommand{\C}{\mathbb{C}}
\newcommand{\bbA}{\mathbb{A}}
\newcommand{\CB}{\mathcal{B}}
\newcommand{\CC}{\mathcal{C}}
\newcommand{\CD}{\mathcal{D}}
\newcommand{\rmC}{\mathrm{C}}
\newcommand{\calX}{\mathcal{X}}

\newcommand{\id}{\operatorname{id}}
\newcommand{\soc}{\operatorname{soc}\nolimits}
\newcommand{\rad}{\operatorname{rad}\nolimits}
\renewcommand{\top}{\operatorname{top}\nolimits}
\newcommand{\Ext}{\operatorname{Ext}\nolimits}
\newcommand{\Tor}{\operatorname{Tor}\nolimits}
\newcommand{\Hom}{\operatorname{Hom}\nolimits}
\newcommand{\End}{\operatorname{End}\nolimits}
\newcommand{\Aut}{\operatorname{Aut}\nolimits}
\newcommand{\op}{\mathrm{op}}
\renewcommand{\Im}{\operatorname{Im}\nolimits}
\newcommand{\Ker}{\operatorname{Ker}\nolimits}
\newcommand{\Cok}{\operatorname{Cok}\nolimits}

\DeclareMathOperator{\moduleCategory}{\mathsf{mod}} \renewcommand{\mod}{\moduleCategory}
\DeclareMathOperator{\Mod}{\mathsf{Mod}}
\DeclareMathOperator{\proj}{\mathsf{proj}}
\DeclareMathOperator{\inj}{\mathsf{inj}}
\DeclareMathOperator{\ind}{\mathsf{ind}}

\DeclareMathOperator{\Char}{char}
\newcommand{\triv}{\operatorname{triv}\nolimits}
\newcommand{\sgn}{\operatorname{sgn}\nolimits}
\newcommand{\Mat}{\operatorname{Mat}\nolimits}
\newcommand{\GL}{\operatorname{GL}\nolimits}
\DeclareMathOperator{\Sn}{\mathfrak{S}}
\newcommand{\Res}{\operatorname{Res}\nolimits}
\newcommand{\Ind}{\operatorname{Ind}\nolimits}
\newcommand{\Symm}{\operatorname{Sym}\nolimits}
\newcommand{\Stab}{\operatorname{Stab}\nolimits}
\newcommand{\Tr}{\operatorname{Tr}\nolimits}
\newcommand{\bfm}{\mathbf{m}}
\newcommand{\bfT}{\mathbf{T}}

\newcommand{\Ann}{\operatorname{Ann}\nolimits}
\newcommand{\linA}{\vec{\mathbb{A}}}
\newcommand{\pdim}{\operatorname{pdim}}
\newcommand{\idim}{\operatorname{idim}}
\newcommand{\gldim}{\operatorname{gldim}\nolimits}
\newcommand{\Rej}{\operatorname{Rej}\nolimits}
\newcommand{\std}{\Delta}
\newcommand{\costd}{\nabla}
\newcommand{\filt}{\mathcal{F}}
\newcommand{\add}{\operatorname{add}}
\newcommand{\dimvec}{\mathbf{dim}}
\newcommand{\nunlhd}{\ntrianglelefteq}
\newcommand{\nlhd}{\ntriangleleft}
\newcommand{\nunrhd}{\ntrianglerighteq}
\newcommand{\nrhd}{\ntriangleright}
\newcommand{\LL}{\mathrm{LL}}

\newcommand{\Ab}{\operatorname{\mathcal{A}b}}

\newcommand{\comm}[1]{{\color[rgb]{0.1,0.9,0.1}#1}}
\newcommand{\old}[1]{{\color{red}#1}}
\newcommand{\new}[1]{{\color{blue}#1}}
\definecolor{darkblue}{rgb}{0,0,0.7}
\newcommand{\defn}[1]{\textsl{\color{darkblue} #1}}

\makeindex

\begin{document}
\begin{center}
\textsc{{\LARGE Topics in Mathematical Science}\\
Spring 2026}\\
\textsc{\Large Introduction to Koszul Duality}

\textsc{Aaron Chan}
\end{center}

Last update: \today

%\tableofcontents

\section*{Convention}

Throughout the course, $\Bbbk$ is a field.
All rings are unital and associative.
All modules are right modules unless otherwise specified.
When we discuss dimensions, we mean dimensions as $\Bbbk$-vector spaces.

\section{Lecture 1: Graded rings and modules}

\subsection*{Algebras}

We begin with the basic object of the course.

\begin{definition}
A \defn{$\Bbbk$-algebra} is a ring $A$ together with a ring homomorphism
\[
\Bbbk \longrightarrow Z(A),
\]
where $Z(A)=\{a\in A\mid ab=ba\text{ for all }b\in A\}$ is the center of $A$.
Equivalently, $A$ is a $\Bbbk$-vector space with an associative multiplication
\[
A\times A\to A,\qquad (a,b)\mapsto ab,
\]
which is $\Bbbk$-bilinear and has an identity element $1_A$.
\end{definition}

The equivalence in the definition is useful.
The ring point of view remembers multiplication; the vector space point of view makes clear that multiplication distributes over linear combinations.

\begin{example}
\begin{enumerate}
\item The field $\Bbbk$ is a $\Bbbk$-algebra.
\item The polynomial ring $\Bbbk[x]$ is a $\Bbbk$-algebra.
\item The matrix ring $\Mat_n(\Bbbk)$ is a $\Bbbk$-algebra.  It is usually non-commutative when $n\geq 2$.
\item If $V$ is a $\Bbbk$-vector space, then $\End_{\Bbbk}(V)$ is a $\Bbbk$-algebra under composition.
\item The quotient $\Bbbk[x]/(x^2)$ is a $\Bbbk$-algebra.  It has basis $\{1,x\}$ and relation $x^2=0$.
\end{enumerate}
\end{example}

\begin{definition}
Let $A$ be a $\Bbbk$-algebra.
A \defn{right $A$-module} is a $\Bbbk$-vector space $M$ together with a map
\[
M\times A\to M,\qquad (m,a)\mapsto ma
\]
such that, for all $m,n\in M$, $a,b\in A$, and $\lambda\in\Bbbk$,
\[
(m+n)a=ma+na,\quad m(a+b)=ma+mb,\quad (\lambda m)a=\lambda(ma),
\]
\[
m(ab)=(ma)b,\qquad m1_A=m.
\]
\end{definition}

Thus a module is a vector space on which elements of $A$ act as linear operators, compatibly with multiplication in $A$.
For instance, a module over $\Bbbk[x]$ is the same thing as a vector space $M$ together with one linear operator $T:M\to M$, where $m x:=T(m)$.

\subsection*{Graded algebras}

Koszul duality is a graded theory.
The grading records a notion of degree, and the main results compare algebraic multiplication with homological degree.

\begin{definition}
A \defn{$\Z_{\geq 0}$-graded $\Bbbk$-algebra} is a $\Bbbk$-algebra $A$ with a direct sum decomposition
\[
A=\bigoplus_{i\geq 0} A_i
\]
as $\Bbbk$-vector spaces such that
\[
A_iA_j\subset A_{i+j}
\]
for all $i,j\geq 0$, and $1_A\in A_0$.
An element $a\in A_i$ is called \defn{homogeneous of degree $i$}.
\end{definition}

Every element $a\in A$ is a finite sum
\[
a=a_0+a_1+\cdots +a_N,\qquad a_i\in A_i.
\]
The summands $a_i$ are called the \defn{homogeneous components} of $a$.

\begin{example}
The polynomial ring $\Bbbk[x]$ is graded by degree:
\[
\Bbbk[x]=\bigoplus_{i\geq 0}\Bbbk x^i.
\]
Here $A_i=\Bbbk x^i$.
More generally,
\[
\Bbbk[x_1,\ldots,x_n]=\bigoplus_{i\geq 0}\Bbbk[x_1,\ldots,x_n]_i,
\]
where $\Bbbk[x_1,\ldots,x_n]_i$ is the span of monomials of total degree $i$.
\end{example}

\begin{example}
Let $V$ be a $\Bbbk$-vector space.
The \defn{tensor algebra} on $V$ is
\[
T(V):=\bigoplus_{i\geq 0}V^{\otimes i}
\]
where $V^{\otimes 0}:=\Bbbk$.
Multiplication is concatenation of tensors:
\[
(v_1\otimes\cdots\otimes v_i)(w_1\otimes\cdots\otimes w_j)
=v_1\otimes\cdots\otimes v_i\otimes w_1\otimes\cdots\otimes w_j.
\]
Thus $T(V)_i=V^{\otimes i}$.
If $V=\Bbbk x\oplus \Bbbk y$, then $T(V)$ is the free non-commutative polynomial algebra $\Bbbk\langle x,y\rangle$.
For example, $xy$ and $yx$ are different basis elements of degree $2$.
\end{example}

\begin{example}
The exterior algebra on an $n$-dimensional vector space $V$ is
\[
\bigwedge V=\bigoplus_{i=0}^{n}\bigwedge^i V.
\]
It is graded by exterior degree.
If $V=\Bbbk x\oplus \Bbbk y$, then
\[
\bigwedge V=\Bbbk\oplus \Bbbk x\oplus \Bbbk y\oplus \Bbbk(x\wedge y),
\]
with $x^2=y^2=0$ and $xy=-yx$.
\end{example}

\begin{definition}
Let $A=\bigoplus_{i\geq 0}A_i$ and $B=\bigoplus_{i\geq 0}B_i$ be graded $\Bbbk$-algebras.
A \defn{homomorphism of graded algebras} $f:A\to B$ is a $\Bbbk$-algebra homomorphism such that
\[
f(A_i)\subset B_i
\]
for all $i\geq 0$.
\end{definition}

\begin{definition}
A graded algebra $A=\bigoplus_{i\geq 0}A_i$ is called \defn{connected} if $A_0=\Bbbk$.
It is called \defn{locally finite} if $\dim_\Bbbk A_i<\infty$ for all $i$.
\end{definition}

Most examples in the first part of the course will be connected and locally finite.
Connectedness means that there is only one copy of the ground field in degree $0$.
Local finiteness means that each degree can be handled by finite-dimensional linear algebra.

\begin{example}
\begin{enumerate}
\item $\Bbbk[x_1,\ldots,x_n]$ is connected and locally finite.
\item $T(V)$ is connected and locally finite if $\dim_\Bbbk V<\infty$.
\item $\Bbbk[x]$ with $\deg x=2$ is also a graded algebra, but it is not generated in degree $1$.
\item The Laurent polynomial ring $\Bbbk[x,x^{-1}]$ is naturally $\Z$-graded, but not $\Z_{\geq 0}$-graded in the above sense.
\end{enumerate}
\end{example}

\subsection*{Graded modules}

\begin{definition}
Let $A=\bigoplus_{i\geq 0}A_i$ be a graded algebra.
A \defn{graded right $A$-module} is a right $A$-module $M$ with a direct sum decomposition
\[
M=\bigoplus_{j\in\Z}M_j
\]
such that
\[
M_jA_i\subset M_{i+j}
\]
for all $i\geq 0$ and all $j\in\Z$.
An element of $M_j$ is called \defn{homogeneous of degree $j$}.
\end{definition}

We allow negative degrees for modules because shifts are unavoidable.
Even if an algebra only has non-negative degrees, its modules may naturally be shifted left or right.

\begin{example}
The algebra $A$ itself is a graded right $A$-module:
\[
A=\bigoplus_{i\geq 0}A_i.
\]
More generally, $A^{\oplus r}$ is a graded free module with each basis vector placed in degree $0$.
\end{example}

\begin{definition}
Let $M$ be a graded $A$-module and let $d\in\Z$.
The \defn{shift} $M(d)$ is the graded $A$-module whose underlying $A$-module is $M$, but whose grading is
\[
M(d)_j=M_{j+d}.
\]
\end{definition}

This convention means that $A(-d)$ has a generator in degree $d$.
Indeed, $1\in A_0=A(-d)_d$.

\begin{example}
Let $A=\Bbbk[x]$ with $\deg x=1$.
The graded module $A(-2)$ has basis
\[
1,x,x^2,\ldots
\]
in degrees
\[
2,3,4,\ldots
\]
respectively.
Thus multiplication by $x$ still raises degree by $1$.
\end{example}

\begin{definition}
Let $M,N$ be graded $A$-modules.
A homomorphism $f:M\to N$ is \defn{homogeneous of degree $d$} if
\[
f(M_j)\subset N_{j+d}
\]
for all $j\in\Z$.
A \defn{homomorphism of graded modules} means a homogeneous $A$-module homomorphism of degree $0$.
\end{definition}

\begin{example}
For $A=\Bbbk[x]$, multiplication by $x^r$ gives a homogeneous map
\[
A\to A
\]
of degree $r$.
The same multiplication can be regarded as a degree $0$ map
\[
A(-r)\to A.
\]
This is why shifts appear constantly in graded free resolutions.
\end{example}

\subsection*{Homogeneous ideals and quotients}

\begin{definition}
Let $A=\bigoplus_{i\geq 0}A_i$ be a graded algebra.
A right ideal $I\subset A$ is \defn{homogeneous} if
\[
I=\bigoplus_{i\geq 0}(I\cap A_i).
\]
Equivalently, whenever $a\in I$ and $a=a_0+\cdots+a_N$ is its decomposition into homogeneous components, each $a_i$ also belongs to $I$.
\end{definition}

For quotient algebras we need two-sided ideals.
If $I$ is a homogeneous two-sided ideal, then
\[
A/I=\bigoplus_{i\geq 0} A_i/(I\cap A_i)
\]
is again a graded algebra.

\begin{example}
In $\Bbbk[x,y]$ with the usual grading, the ideal
\[
I=(x^2,xy,y^3)
\]
is homogeneous.
The quotient $A=\Bbbk[x,y]/I$ is graded.
Its degree $0$ part is $\Bbbk$, its degree $1$ part has basis $x,y$, and its degree $2$ part has basis $y^2$.
All degrees $\geq 3$ vanish.
\end{example}

\begin{example}
The ideal $(x-1)\subset \Bbbk[x]$ is not homogeneous for the usual grading.
Indeed $x-1$ is in the ideal, but its homogeneous components $x$ and $-1$ are not in the ideal.
The quotient $\Bbbk[x]/(x-1)\cong \Bbbk$ should not be regarded as a graded quotient of $\Bbbk[x]$ by this ideal.
\end{example}

\begin{definition}
Let $A=\bigoplus_{i\geq 0}A_i$ be a connected graded algebra.
The ideal
\[
A_+:=\bigoplus_{i>0}A_i
\]
is called the \defn{augmentation ideal}.
The quotient
\[
\Bbbk_A:=A/A_+
\]
is called the \defn{trivial module}.
\end{definition}

The trivial module is the simplest graded module over $A$.
Every positive degree element acts by zero.
Koszulness will later be a statement about the shape of the minimal graded free resolution of this module.

\subsection*{Hilbert series}

\begin{definition}
Let $M=\bigoplus_{j\in\Z}M_j$ be a graded vector space with $\dim_\Bbbk M_j<\infty$ for all $j$ and $M_j=0$ for all $j\ll 0$.
The \defn{Hilbert series} of $M$ is the formal Laurent series
\[
h_M(t):=\sum_{j\in\Z}(\dim_\Bbbk M_j)t^j.
\]
For a locally finite graded algebra $A$, we write $h_A(t)$ for its Hilbert series as a graded vector space.
\end{definition}

The Hilbert series records the size of each degree.
It forgets multiplication, but it is still strong enough to give numerical tests for Koszulness later.

\begin{lemma}
Let $M,N$ be locally finite graded vector spaces.
\begin{enumerate}
\item If $M\cong N$ as graded vector spaces, then $h_M(t)=h_N(t)$.
\item If $0\to L\to M\to N\to 0$ is an exact sequence of graded vector spaces with degree $0$ maps, then
\[
h_M(t)=h_L(t)+h_N(t).
\]
\item For any shift $d\in\Z$,
\[
h_{M(d)}(t)=t^{-d}h_M(t).
\]
\end{enumerate}
\end{lemma}

\begin{proof}
All three statements follow by taking the degree $j$ part.
For example, $M(d)_j=M_{j+d}$, so
\[
h_{M(d)}(t)=\sum_j \dim_\Bbbk(M_{j+d})t^j=t^{-d}\sum_i \dim_\Bbbk(M_i)t^i.
\]
\end{proof}

\begin{example}
For $A=\Bbbk[x]$ with $\deg x=1$, we have $\dim_\Bbbk A_i=1$ for all $i\geq 0$.
Hence
\[
h_A(t)=1+t+t^2+\cdots=\frac{1}{1-t}.
\]
For $\Bbbk[x_1,\ldots,x_n]$ with all variables in degree $1$,
\[
h_{\Bbbk[x_1,\ldots,x_n]}(t)=\frac{1}{(1-t)^n}.
\]
Indeed, this follows by multiplying the $n$ geometric series, one for each variable.
\end{example}

\begin{example}
Let $V$ be $n$-dimensional.
Then $\dim_\Bbbk T(V)_i=n^i$, so
\[
h_{T(V)}(t)=1+nt+n^2t^2+\cdots=\frac{1}{1-nt}.
\]
On the other hand,
\[
h_{\bigwedge V}(t)=\sum_{i=0}^{n}\binom{n}{i}t^i=(1+t)^n.
\]
\end{example}

\begin{example}
Let $A=\Bbbk[x,y]/(x^2,xy,y^3)$.
From the basis calculation above,
\[
h_A(t)=1+2t+t^2.
\]
This series already tells us that $A$ is $4$-dimensional as a vector space.
\end{example}

\begin{lemma}
If $V,W$ are locally finite graded vector spaces, and each graded piece of $V\otimes_\Bbbk W$ is finite-dimensional, then
\[
h_{V\otimes W}(t)=h_V(t)h_W(t),
\]
where $(V\otimes W)_n=\bigoplus_{i+j=n}V_i\otimes W_j$.
\end{lemma}

\begin{proof}
The degree $n$ part has dimension
\[
\dim_\Bbbk (V\otimes W)_n=\sum_{i+j=n}\dim_\Bbbk V_i\dim_\Bbbk W_j.
\]
This is exactly the coefficient of $t^n$ in the product $h_V(t)h_W(t)$.
\end{proof}

\subsection*{Graded duals}

\begin{definition}
Let $V=\bigoplus_{i\in\Z}V_i$ be a locally finite graded vector space.
Its \defn{graded dual} is
\[
V^\vee:=\bigoplus_{i\in\Z}\Hom_\Bbbk(V_i,\Bbbk).
\]
We regard $\Hom_\Bbbk(V_i,\Bbbk)$ as sitting in degree $-i$.
\end{definition}

The ordinary dual $\Hom_\Bbbk(V,\Bbbk)$ is a direct product of the duals of the graded pieces.
The graded dual is the direct sum.
For finite-dimensional graded spaces there is no difference as vector spaces, but for infinite-dimensional examples such as $\Bbbk[x]$ the distinction matters.

\begin{example}
Let $A=\Bbbk[x]$ with $\deg x=1$.
The graded dual $A^\vee$ has one-dimensional piece in each degree $0,-1,-2,\ldots$ and no other degrees.
The functional dual to $x^i$ has degree $-i$.
\end{example}

\subsection*{Left, right, and two-sided duals}

The previous definition is enough for connected algebras, where $A_0=\Bbbk$.
However, many natural Koszul algebras have degree zero part a semisimple algebra $S$ rather than just $\Bbbk$.
For example, the path algebra of a quiver has
\[
S=\bigoplus_{i\in Q_0}\Bbbk e_i
\]
in degree $0$.
In that setting one must keep track of left and right actions.

\begin{definition}
Let $S$ be a finite-dimensional semisimple $\Bbbk$-algebra.
\begin{enumerate}
\item If ${}_S V$ is a finite-dimensional left $S$-module, define its \defn{right $S$-dual}
\[
V^*:=\Hom_S(V,S).
\]
It is a right $S$-module by
\[
(fs)(v):=f(v)s.
\]

\item If $W_S$ is a finite-dimensional right $S$-module, define its \defn{left $S$-dual}
\[
{}^*W:=\Hom_{S^\op}(W,S).
\]
It is a left $S$-module by
\[
(sg)(w):=sg(w).
\]
\end{enumerate}
\end{definition}

Thus duality reverses sides.
This is already familiar over a field: if $M$ is a right $A$-module, then $\Hom_\Bbbk(M,\Bbbk)$ is naturally a left $A$-module by
\[
(af)(m):=f(ma).
\]
The same formula explains the general case.

\begin{lemma}
Let $S$ be finite-dimensional semisimple.
For finite-dimensional modules, the evaluation maps
\[
V\to {}^*(V^*),\qquad W\to ({}^*W)^*
\]
are isomorphisms.
\end{lemma}

\begin{proof}
Since $S$ is semisimple, every finite-dimensional $S$-module is a direct summand of a finite direct sum of copies of $S$.
The statement is immediate for $S$ itself, and then follows for finite direct sums and direct summands.
\end{proof}

\begin{definition}
An \defn{$S$-bimodule} is a vector space $V$ which is both a left and right $S$-module, with
\[
(sv)t=s(vt)
\]
for all $s,t\in S$ and $v\in V$.
Equivalently, it is a left module over the \defn{enveloping algebra}
\[
S^{\mathrm{en}}:=S\otimes_\Bbbk S^\op.
\]
\end{definition}

The bimodule viewpoint is the correct two-sided language.
If $V$ is an $S$-bimodule, then both $V^*$ and ${}^*V$ are again $S$-bimodules, but the missing action must be defined carefully:
\[
(sf)(v):=f(vs)\qquad \text{for }f\in V^*,
\]
and
\[
(gt)(v):=g(tv)\qquad \text{for }g\in {}^*V.
\]
Together with the actions in the previous definition, these formulas make $V^*$ and ${}^*V$ into bimodules.

\begin{example}
Take $S=\Bbbk^r$ with primitive idempotents $e_1,\ldots,e_r$.
An $S$-bimodule decomposes as
\[
V=\bigoplus_{i,j}e_iVe_j.
\]
The summand $e_iVe_j$ should be thought of as arrows whose left end is $i$ and right end is $j$.
Duality reverses the two sides:
\[
e_jV^*e_i\cong \Hom_\Bbbk(e_iVe_j,\Bbbk).
\]
This is the linear algebra reason that dual quivers reverse arrows.
\end{example}

\begin{example}
If $S=\Bbbk^r$ with primitive idempotents $e_1,\ldots,e_r$, then an $S$-module is the same as a direct sum of vector spaces
\[
V=Ve_1\oplus\cdots\oplus Ve_r.
\]
Duality over $S$ dualizes each component separately.
This is the linear algebra shadow of what happens for quivers with several vertices.
\end{example}

Tensor products also reverse order under duality.
For finite-dimensional $S$-bimodules there are natural isomorphisms
\[
W^*\otimes_S V^*\cong (V\otimes_S W)^*,
\qquad
{}^*W\otimes_S {}^*V\cong {}^*(V\otimes_S W).
\]
The reversal is not cosmetic: a functional on $V\otimes_S W$ first sees an element of $V$, then an element of $W$, while the dual tensor product records the corresponding operations in the opposite order.

\begin{example}
When $S=\Bbbk$, this reduces to the usual isomorphism
\[
W^*\otimes_\Bbbk V^*\cong (V\otimes_\Bbbk W)^*
\]
for finite-dimensional vector spaces.
The order reversal is usually hidden because vector spaces have no left/right distinction.
\end{example}

\begin{definition}
Let $V$ be a finite-dimensional $S$-bimodule and let $U\subset V$ be a sub-bimodule.
Define the \defn{right orthogonal complement}
\[
U^\perp:=\{f\in V^*\mid f(U)=0\},
\]
and the \defn{left orthogonal complement}
\[
{}^\perp U:=\{g\in {}^*V\mid g(U)=0\}.
\]
Both are $S$-sub-bimodules.
\end{definition}

This is the notation needed for quadratic duality.
If
\[
A=T_S(V)/(R),\qquad R\subset V\otimes_S V,
\]
then the \defn{left quadratic dual} in the sense of BGS is
\[
A^!:=T_S(V^*)/(R^\perp),
\]
where $R^\perp\subset V^*\otimes_S V^*\cong (V\otimes_S V)^*$.
There is also a \defn{right quadratic dual}
\[
{}^!A:=T_S({}^*V)/({}^\perp R).
\]
When $S=\Bbbk$, the two constructions identify canonically with the usual quadratic dual.
For general semisimple $S$, distinguishing them prevents left modules, right modules, and bimodules from being silently interchanged.

\pagebreak

\section{Lecture 2: Tensor and quadratic algebras}

\subsection*{Generated algebras}

Let $A=\bigoplus_{i\geq 0}A_i$ be a connected graded algebra.
The multiplication maps give linear maps
\[
A_1^{\otimes i}\to A_i,\qquad a_1\otimes\cdots\otimes a_i\mapsto a_1\cdots a_i.
\]
If these maps are surjective for all $i\geq 0$, then every homogeneous element of positive degree can be written as a linear combination of products of degree $1$ elements.

\begin{definition}
A connected graded algebra $A$ is \defn{generated in degree $1$} if the smallest graded subalgebra of $A$ containing $A_0=\Bbbk$ and $A_1$ is all of $A$.
Equivalently, the maps
\[
A_1^{\otimes i}\to A_i
\]
are surjective for all $i\geq 0$.
\end{definition}

\begin{example}
\begin{enumerate}
\item $\Bbbk[x_1,\ldots,x_n]$ with $\deg x_i=1$ is generated in degree $1$.
\item $T(V)$ is generated in degree $1$, since every tensor is a product of elements of $V=T(V)_1$.
\item $\bigwedge V$ is generated in degree $1$.
\item $\Bbbk[x]$ with $\deg x=2$ is not generated in degree $1$, since $A_1=0$ but $A_2=\Bbbk x\neq 0$.
\end{enumerate}
\end{example}

The main construction today explains why tensor algebras are the natural starting point.
Every algebra generated in degree $1$ is a quotient of a tensor algebra.

\subsection*{Tensor algebras and the universal property}

\begin{definition}
Let $V$ be a $\Bbbk$-vector space.
The \defn{tensor algebra} of $V$ is
\[
T(V):=\bigoplus_{i\geq 0}V^{\otimes i},
\]
with $V^{\otimes 0}=\Bbbk$ and multiplication given by concatenation.
\end{definition}

Thus $T(V)$ consists of finite linear combinations of words in elements of $V$.
It is the algebra with no relations among the generators except those forced by linearity and associativity.

\begin{example}
If $V=\Bbbk x$, then
\[
T(V)\cong \Bbbk[x].
\]
There is only one letter, so non-commutativity is invisible.
If $V=\Bbbk x\oplus \Bbbk y$, then
\[
T(V)\cong \Bbbk\langle x,y\rangle,
\]
with basis all words in $x$ and $y$:
\[
1,\quad x,y,\quad x^2,xy,yx,y^2,\quad x^3,x^2y,xyx,\ldots.
\]
\end{example}

\begin{proposition}[Universal property of the tensor algebra]
Let $B$ be a $\Bbbk$-algebra and let $f:V\to B$ be a $\Bbbk$-linear map.
There exists a unique $\Bbbk$-algebra homomorphism
\[
\widetilde f:T(V)\to B
\]
such that $\widetilde f(v)=f(v)$ for all $v\in V$.
\end{proposition}

\begin{proof}
Define
\[
\widetilde f(v_1\otimes\cdots\otimes v_i)=f(v_1)\cdots f(v_i)
\]
on pure tensors and extend linearly.
This gives a homomorphism of algebras because concatenation of tensors is sent to multiplication in $B$.
The formula is forced on every product of elements of $V$, so the homomorphism is unique.
\end{proof}

\begin{corollary}
Let $A$ be a connected graded algebra generated in degree $1$.
Then there is a surjective homomorphism of graded algebras
\[
T(A_1)\onto A.
\]
\end{corollary}

\begin{proof}
Apply the universal property to the inclusion $A_1\into A$.
Since $A$ is generated in degree $1$, the induced map is surjective.
\end{proof}

\begin{definition}
Let $A$ be a connected graded algebra generated in degree $1$.
Choose $V=A_1$ and write
\[
A\cong T(V)/I
\]
for a homogeneous two-sided ideal $I\subset T(V)$.
The ideal $I$ is called the \defn{ideal of relations}.
\end{definition}

The vector space $V$ gives the generators.
The ideal $I$ tells us which linear combinations of words are declared to be zero.
This is the same principle as writing
\[
\Bbbk[x,y]\cong \Bbbk\langle x,y\rangle/(xy-yx).
\]

\subsection*{Relations in degree 2}

Since $T(V)$ is graded, every homogeneous relation has a degree.
Relations in degree $2$ are linear combinations of words of length $2$, hence elements of $V\otimes V$.

\begin{definition}
Let $V$ be a finite-dimensional $\Bbbk$-vector space and let
\[
R\subset V\otimes V
\]
be a subspace.
The algebra
\[
A=T(V)/(R)
\]
is called a \defn{quadratic algebra}.
Here $(R)$ denotes the two-sided ideal of $T(V)$ generated by $R$.
\end{definition}

In other words, a quadratic algebra is an algebra generated by degree $1$ elements whose defining relations are homogeneous relations of degree $2$.
The word ``quadratic'' refers to degree, not to commutativity.

\begin{example}
Let $V=\Bbbk x\oplus \Bbbk y$.
The commutative polynomial ring is
\[
\Bbbk[x,y]\cong T(V)/(x\otimes y-y\otimes x).
\]
We usually write this as
\[
\Bbbk[x,y]\cong \Bbbk\langle x,y\rangle/(xy-yx).
\]
Thus the polynomial ring is quadratic.
\end{example}

\begin{example}
Let $V$ have basis $x_1,\ldots,x_n$.
The polynomial ring is
\[
\Bbbk[x_1,\ldots,x_n]\cong
T(V)\Big/
\left(x_ix_j-x_jx_i\mid 1\leq i<j\leq n\right).
\]
There are $\binom n2$ independent quadratic commutativity relations.
\end{example}

\begin{example}
Let $V$ have basis $x_1,\ldots,x_n$.
The exterior algebra is
\[
\bigwedge V\cong
T(V)\Big/
\left(x_i^2,\ x_ix_j+x_jx_i\mid 1\leq i<j\leq n\right),
\]
provided $\Char\Bbbk\neq 2$.
If $\Char\Bbbk=2$, the signs require more care; one should define the exterior algebra by the alternating relation $v\otimes v=0$ for all $v\in V$.
\end{example}

\begin{example}
Let
\[
A=\Bbbk\langle x,y\rangle/(x^2,xy).
\]
Then $A$ is quadratic.
A basis of $A$ is given by words in $x,y$ which do not contain $x^2$ or $xy$ as consecutive subwords.
Thus after an $x$ appears, no letter may follow it.
The first few basis elements are
\[
1,\quad x,y,\quad yx,y^2,\quad y^2x,y^3,\quad \ldots.
\]
Hence
\[
h_A(t)=1+2t+2t^2+2t^3+\cdots=1+\frac{2t}{1-t}.
\]
\end{example}

\begin{example}
The algebra
\[
\Bbbk[x]/(x^3)
\]
with $\deg x=1$ is generated in degree $1$, but it is not quadratic.
Its defining relation has degree $3$.
This algebra is important in ring theory, but it is not in the basic class governed by quadratic duality.
\end{example}

\begin{example}
The algebra
\[
\Bbbk[x]/(x^2)
\]
with $\deg x=1$ is quadratic:
\[
\Bbbk[x]/(x^2)\cong T(\Bbbk x)/(x\otimes x).
\]
As a graded vector space it has basis $1,x$, so
\[
h_A(t)=1+t.
\]
\end{example}

\subsection*{The degree 2 relations of a quotient}

Let $A=T(V)/I$ be a connected graded algebra generated in degree $1$.
The degree $2$ part of $T(V)$ is $V\otimes V$.
The degree $2$ part of the ideal $I$ is
\[
I_2:=I\cap (V\otimes V).
\]
The quotient map gives
\[
A_2\cong (V\otimes V)/I_2.
\]
Thus $I_2$ is exactly the space of quadratic relations visible in $A$.

\begin{definition}
Let $A=T(V)/I$ be generated in degree $1$.
The \defn{quadratic part} of $A$ is
\[
qA:=T(V)/(I_2).
\]
There is a natural surjective homomorphism
\[
qA\onto A.
\]
\end{definition}

The algebra $qA$ keeps the same generators and the same degree $2$ relations as $A$, but forgets all relations of degree $3$ or higher.
Thus $A$ is quadratic precisely when $qA\to A$ is an isomorphism.

\begin{example}
Let $A=\Bbbk[x]/(x^3)$ with $\deg x=1$.
Then $V=\Bbbk x$ and $I=(x^3)$.
Since $I_2=0$, we have
\[
qA=T(V)\cong \Bbbk[x].
\]
The natural map $qA\to A$ is the quotient map $\Bbbk[x]\to \Bbbk[x]/(x^3)$, not an isomorphism.
\end{example}

\begin{example}
Let
\[
A=\Bbbk\langle x,y\rangle/(xy-yx,\ x^3).
\]
Then $A$ is generated in degree $1$ but not quadratic.
Its quadratic part is
\[
qA=\Bbbk\langle x,y\rangle/(xy-yx)\cong \Bbbk[x,y].
\]
The cubic relation $x^3=0$ is invisible to $qA$.
\end{example}

\subsection*{Why quadratic algebras are the right class}

Koszul duality begins with the observation that degree $2$ relations have a linear dual.
If $A=T(V)/(R)$ with $R\subset V\otimes V$, then the dual vector space $V^*$ gives dual generators, while the orthogonal subspace
\[
R^\perp\subset V^*\otimes V^*
\]
will give the dual quadratic relations.
This construction will be made precise next lecture.

For now, note why degree $2$ is special.
The dual of $V\otimes V$ is naturally $V^*\otimes V^*$ when $V$ is finite-dimensional.
Thus a subspace of quadratic relations has an orthogonal complement which is again a space of quadratic relations.
There is no similarly simple operation that turns arbitrary cubic and quartic relations into relations of the same kind without adding more structure.

\begin{example}
For the polynomial ring
\[
\Bbbk[x,y]\cong \Bbbk\langle x,y\rangle/(xy-yx),
\]
the space of quadratic relations is one-dimensional, spanned by $xy-yx$.
The quadratic dual will have two generators $x^*,y^*$ and relations corresponding to the orthogonal complement of this line in
\[
V^*\otimes V^*.
\]
This produces the exterior algebra on two generators.
The full computation will be done in Lecture 3.
\end{example}

\begin{example}
For the exterior algebra
\[
\bigwedge(\Bbbk x\oplus\Bbbk y)
\cong \Bbbk\langle x,y\rangle/(x^2,\ y^2,\ xy+yx),
\]
the quadratic dual will be the polynomial ring in two variables.
This is the first sign of the duality between symmetric and exterior algebras.
\end{example}

\subsection*{First checks and non-examples}

\begin{proposition}
Let $A=T(V)/(R)$ be quadratic, where $V$ is finite-dimensional and $R\subset V\otimes V$.
Then $A$ is connected, locally finite, and generated in degree $1$.
\end{proposition}

\begin{proof}
The tensor algebra $T(V)$ is connected and generated in degree $1$.
The ideal $(R)$ is homogeneous because it is generated by homogeneous elements of degree $2$.
Therefore the quotient inherits a connected grading and is generated by the image of $V$.
Since $V$ is finite-dimensional, each $T(V)_i=V^{\otimes i}$ is finite-dimensional, and hence each quotient $A_i$ is finite-dimensional.
\end{proof}

\begin{remark}
The converse is false.
A connected locally finite algebra generated in degree $1$ need not be quadratic.
The example $\Bbbk[x]/(x^3)$ is the smallest warning.
\end{remark}

\begin{exercise}
Let $A=\Bbbk\langle x,y\rangle/(x^2,y^2,xy)$.
Find a basis for $A_i$ for $i=0,1,2,3$, and compute the Hilbert series.
\end{exercise}

\begin{exercise}
Let $A=\Bbbk[x,y]/(x^2,xy,y^2)$ with the usual grading.
Show that $A$ is quadratic and compute $h_A(t)$.
\end{exercise}

\begin{exercise}
Let $V$ be $n$-dimensional.
Show directly that
\[
\dim_\Bbbk \Symm^2(V)=\frac{n(n+1)}{2},\qquad
\dim_\Bbbk \bigwedge^2 V=\frac{n(n-1)}{2}.
\]
Interpret these dimensions as the number of independent quadratic monomials in the commutative and anti-commutative cases.
\end{exercise}

\end{document}
